\section{Limitations}
\label{sec:limitations}

\paragraph{Model and trigger coverage.}
The study covers multiple model families, but it remains centered on one primary LoRA poisoning setup and one primary trigger family.
Additional sweeps vary trigger syntax and poison ratio, but their results are mixed, so we do not treat them as broad trigger-level generalization.
Future work should broaden trigger syntax, poison ratio, target behavior, adapter placement, and fine-tuning recipe under the same controlled four-way protocol.
The Llama result suggests that model architecture can substantially change the attainable safety-utility frontier.

\paragraph{Dataset disclosure.}
The formal four-way set is human-reviewed and identified by a stable SHA-256 hash, and the public release includes a dataset card, schema, annotation guidelines, sanitized examples, and a sampling manifest.
We do not publish the raw JSONL file because it contains harmful requests, exact trigger interventions, and potentially unsafe model generations.
This protects against misuse but means that full raw-prompt metric reproduction requires controlled access or substitution with an approved harmful-behavior benchmark and synthetic trigger interface.

\paragraph{Evaluation protocol and judge robustness.}
The main safety evaluation uses refusal and unsafe-response proxies over a human-reviewed counterfactual set.
This protocol is useful for controlled large-scale comparison, but LLM judges can be brittle for borderline generations, partial compliance, and ambiguous refusals.
Appendix~\ref{app:judge-audit} reports a 200-example secondary local-judge audit; agreement is high overall, while borderline TH partial-compliance cases remain the main source of disagreement.
The Llama MMLU values are reported from a separate utility run, which should be kept explicit when comparing utility across model families.

\paragraph{Statistical coverage.}
The main rows are formal 1,000-example evaluations, and the matched random-pruning baseline is repeated over ten seeds.
The Qwen27B margins are large under both binomial-rate uncertainty and random-budget variability.
The remaining statistical caveat is that most non-random \sac{} routes are single selected adapters; their cross-model differences should therefore be read as frontier evidence rather than as estimates of training-seed variance.

\paragraph{Incomplete operator space.}
\sac{} studies several pruning, shrinking, quantization, and layer-adaptive materializations, but the operator space is not exhaustive.
Additional operators, joint optimization of gates and materialization, or calibration with stronger safety judges may improve difficult regimes such as Llama.

\paragraph{Defense-baseline scope.}
The main baselines establish that the effect is not caused by compression pressure alone: behavior-aware component selection is qualitatively different from random pruning, magnitude-energy pruning, low-singular-value pruning, and uniform quantization.
Additional post-hoc defense controls bound nearby alternatives: representation editing leaves the attack active, refusal tuning becomes a near-universal refusal endpoint, and SASP-LoRA pruning is non-degenerate but substantially weaker than \sac{} on Qwen27B TH suppression.
These results compare \sac{} to nearby post-hoc alternatives rather than to every possible backdoor-removal or model-editing defense.

\paragraph{Over-refusal.}
The Qwen4B results show a real, non-degenerate refusal-suppression frontier.
There are normal-refusal points, such as \method{SAC-alpha-60} with TB/B 0.083/0.081, and mid-frontier points such as same-gate soft shrink and \method{SAC-alpha-70}.
Qwen4B additional rows also include cleaner low-TH settings, such as rank-64 and target-all-linear adapters, so Qwen4B should not be reduced to a simple all-refusal pattern.
The caveat is slope and trigger sensitivity: stronger attack-suppression points often incur more triggered-benign refusal than on Qwen27B, and the most aggressive endpoints remain deployment-risky.
This can be a real deployment cost: in a deployed system, trigger-conditioned benign refusal can behave like a denial-of-service channel.
This makes over-refusal an operating cost, not only a metric side effect.

\paragraph{Predicting adapter geometry.}
The results do not identify a universal compression setting for all adapters and base models.
Understanding which architectural or adapter properties predict the attainable frontier remains an open problem.

\paragraph{Mechanism analysis.}
The gate analysis provides initial localization evidence: selected components are structured by layer/projection, the Qwen27B budget curve follows the same ranking, and probe-refit gates overlap the selected gate more than random gates do.
This evidence is mechanistic but not a complete causal geometry of the adapter.
A stronger account would directly intervene on top-scored, bottom-scored, and matched-random components across budgets and architectures to test whether the selected directions separate trigger-supporting behavior from benign or utility-supporting behavior.
